\documentclass{beamer}
\usepackage[utf8]{inputenc}

\usetheme{Madrid}
\usecolortheme{default}
\usepackage{amsmath,amssymb,amsfonts,amsthm}
\usepackage{txfonts}
\usepackage{tkz-euclide}
\usepackage{listings}
\usepackage{adjustbox}
\usepackage{array}
\usepackage{tabularx}
\usepackage{gvv}
\usepackage{multicol}   % <-- Add this
\usepackage{enumitem}

\usepackage{lmodern}
\usepackage{circuitikz}
\usepackage{tikz}
\usepackage{graphicx}
\usepackage{mathtools}
\setbeamertemplate{page number in head/foot}[totalframenumber]

\usepackage{tcolorbox}
\tcbuselibrary{minted,breakable,xparse,skins}



\definecolor{bg}{gray}{0.95}
\DeclareTCBListing{mintedbox}{O{}m!O{}}{%
  breakable=true,
  listing engine=minted,
  listing only,
  minted language=#2,
  minted style=default,
  minted options={%
    linenos,
    gobble=0,
    breaklines=true,
    breakafter=,,
    fontsize=\small,
    numbersep=8pt,
    #1},
  boxsep=0pt,
  left skip=0pt,
  right skip=0pt,
  left=25pt,
  right=0pt,
  top=3pt,
  bottom=3pt,
  arc=5pt,
  leftrule=0pt,
  rightrule=0pt,
  bottomrule=2pt,
  toprule=2pt,
  colback=bg,
  colframe=orange!70,
  enhanced,
  overlay={%
    \begin{tcbclipinterior}
    \fill[orange!20!white] (frame.south west) rectangle ([xshift=20pt]frame.north west);
    \end{tcbclipinterior}},
  #3,
}
\lstset{
    language=C,
    basicstyle=\ttfamily\small,
    keywordstyle=\color{blue},
    stringstyle=\color{orange},
    commentstyle=\color{green!60!black},
    numbers=left,
    numberstyle=\tiny\color{gray},
    breaklines=true,
    showstringspaces=false,
}


\title 
{12.192}
\date{November 3, 2025}


\author 
{Dhanush Kumar A - AI25BTECH11010}



\begin{document}
\frame{\titlepage}
\begin{frame}{Question}


Consider a discrete-time LTI system with input  
$x[n] = \delta[n] + \delta[n - 1]$  
and impulse response  
$h[n] = \delta[n] - \delta[n - 1]$.  

The output of the system will be given by:
\begin{multicols}{2}

\begin{enumerate}
    \item[(a)] $\delta[n] - \delta[n - 2]$
    \item[(b)] $\delta[n] - \delta[n - 1]$
    \item[(c)] $\delta[n - 1] + \delta[n - 2]$
    \item[(d)] $\delta[n] + \delta[n - 1] + \delta[n - 2]$
\end{enumerate}
\end{multicols}
\end{frame}
\begin{frame}{Solution}


For an LTI system, the output is obtained by convolution:
\begin{align}
y[n] &= x[n] * h[n].
\end{align}

Given input and impulse response:
\begin{align}
x[n] &= \delta[n] + \delta[n-1], \\
h[n] &= \delta[n] - \delta[n-1].
\end{align}

Representing them as vectors:
\begin{align}
\vec{x} &= 
\myvec{x[0] \\ x[1]}
= 
\myvec{1 \\ 1}, &
\vec{h} &= 
\myvec{h[0] \\ h[1]}
= 
\myvec{1 \\ -1}.
\end{align}

Determine output length:
\begin{align}
N_y &= N_x + N_h - 1 = 2 + 2 - 1 = 3.
\end{align}
\end{frame}
\begin{frame}{Solution}


Form the convolution (Toeplitz) matrix using $\vec{h}$:
\begin{align}
\vec{H} &= 
\myvec{
1 & 0 \\
-1 & 1 \\
0 & -1
}.
\end{align}

Compute the output vector:
\begin{align}
\vec{y} &= \vec{H}\vec{x} \nonumber \\
&= 
\myvec{
1 & 0 \\
-1 & 1 \\
0 & -1
}
\myvec{1 \\ 1} \nonumber \\
&= 
\myvec{
1(1) + 0(1) \\
-1(1) + 1(1) \\
0(1) - 1(1)
} =
\myvec{1 \\ 0 \\ -1}.
\end{align}
\end{frame}
\begin{frame}{Solution}

Thus, the output samples are:
\begin{align}
y[0] &= 1, &
y[1] &= 0, &
y[2] &= -1.
\end{align}

Hence, the output signal is:
\begin{align}
y[n] &= \delta[n] - \delta[n-2].
\end{align}

Final answer:
\begin{align}
\boxed{y[n] = \delta[n] - \delta[n-2] \quad \text{(Option (a))}}
\end{align}

\end{frame}

\end{document}


